% Nouveau EGA V (cote 25): archive cote 25, site page 166 (PDF page 167).
% typescript. Draft compared: none.
% AI-transcribed by Claude Opus 5.5 (claude-opus-5-5, adaptive thinking, effort high), Message Batches API, batch ega25b (msgbatch_01A7n5YuXubUPNK1zsepi8RG), 2026-11-13; not yet checked by eye.
\scanpage{25/166}
\margin{//} \hand{nouveau EGA V} \mspage{1}

\margin{(}EGA IV par ?

\uline{Ensemble singulier et supersingulier}.

Pages 14 et 15 des notes précédentes, remplacer partout S par Y. Page
15 ligne 10 le \hand{``}i.e.\hand{''} est ab\unsure{us}if, c'est la deuxième assertion qui est la
plus précise et celle à laquelle on en a.

La généralisation de prop.7 à laquelle il est fait allusion pages 16
et 17 s'obtient à l'aide du diagramme suivant, qui généralise celui de la
page 4 :
\diagram{\begin{tikzcd}
 & & 0 & & \\
 & & \underline{\Omega}^1_{X/Y}\otimes\underline{O}_{V^1} \arrow[u] & & \\
0 \arrow[r] & \underline{N}_{V^1/X} \arrow[r] \arrow[ur, "\hand{\mu}"] & \underline{\Omega}^1_{X/S}\otimes\underline{O}_{V^1} \arrow[r] \arrow[u] & \underline{\Omega}^1_{V^1/S} \arrow[r] & 0 \\
 & & \underline{\Omega}^1_{Y/S}\otimes\underline{O}_{V^1} \arrow[u] \arrow[ur, "\hand{\wr}"'] & & \\
 & & \hand{0} \arrow[u] & &
\end{tikzcd}}

où $\underline{N}$ désigne le faisceau conormal. Tous les faisceaux du diagramme sont
localement libres\inserted{sur $V^1$}, et on applique le lemme de la page 7 à $M=\underline{\Omega}^1_{X/S}\otimes\underline{O}_{V^1}$
et aux deux sous-Modules figurant dans le diagramme précédent. Avec
les notations de ce lemme, on est ramené à prouver que la section a
\xout{a} de \uline{Hom}(L,P) déduit de \inserted{$\mu$:} $P\rightarrow M/Q$ i.e. de \inserted{$\mu$:} $\underline{N}_{V^1/X}\longrightarrow\underline{\Omega}^1_{X/Y}\otimes\underline{O}_{V^1}$ ,
dont le \inserted{pré}schéma des zéros V(a) est celui qui rend \hand{$\mu$} strictement de rang
\hand{$<$} d = rang M/Q , satisfait V(a) = V(D), où D est défini au bas de la
page \over{\unsure{1}}{1}5. Or les inclusions de SS-préschémas lisses
\margin{||}
\[ V^1 \subset V^o \subset X \]
donnent une suite exacte (réf ?)
\[ 0 \rightarrow \underline{N}_{V^o/X}\otimes\underline{O}_{V^1} \rightarrow \underline{N}_{V^1/X} \longrightarrow \underline{N}_{V^1/V^o} \rightarrow 0 \ , \]
d'autre part on a des isomorphismes canoniques
\[ \underline{N}_{V^o/X} \simeq \underline{E}\otimes\underline{O}_{V^1} \ , \quad \underline{N}_{V^1/V^o} \simeq \uline{Hom}(W,K) \ , \]
où W,K sont définis (page 15) comme noyau et conoyau de

\margin{166}
